% ==============================================================================
% MÓDULO CAPÍTULO II: CARACTERÍSTICAS SANITARIAS, CARTERA Y CAPACIDAD RESOLUTIVA
% Archivo: cap_2/2.5_sanitaria_capacidad.tex
% ==============================================================================
\subsection{Características sanitarias, cartera de servicios y capacidad resolutiva institucional}
\label{sec:2.5_sanitaria}

El Centro de Salud Belén se encuentra formalmente categorizado como un establecimiento de salud del primer nivel de atención de nivel \textbf{I-3} (sin internamiento hospitalario), desempeñando un papel neurálgico como cabecera de microred asistencial. Su misión operativa consiste en brindar atención integral ambulatoria, promocional, preventiva y de recuperación básica a la población vulnerable de su ámbito territorial y jurisdiccional.

A continuación se detalla la oferta prestacional y la distribución del talento humano asistencial disponible en el establecimiento:

\begin{table}[H]
\centering
\small
\caption{Cartera de servicios asistenciales y talento humano del Centro de Salud Belén (I-3).}
\label{tab:cartera_servicios}
\vspace{0.2cm}
\begin{tabularx}{\textwidth}{p{3.8cm}Xcc}
\toprule
\textbf{Cartera de Servicios} & \textbf{Actividades Asistenciales Principales} & \textbf{Turno} & \textbf{Personal Asignado} \\
\midrule
Medicina General & Consulta ambulatoria, diagnóstico clínico, prescripción médica y atención de morbilidades agudas y crónicas. & Mañana / Tarde & 04 Médicos Cirujanos \\
\addlinespace
Enfermería (CRED e Inmunizaciones) & Control del crecimiento y desarrollo del niño, tamizaje de anemia, vacunación del esquema regular y consejería nutricional. & Mañana / Tarde & 06 Lic. en Enfermería \\
\addlinespace
Enfermería (TBC y No Transmisibles) & Control de tuberculosis (PCT), administración de tratamiento DOTS, tamizaje de hipertensión arterial y diabetes mellitus. & Mañana & 02 Lic. en Enfermería \\
\addlinespace
Obstetricia (Salud Reproductiva) & Control prenatal, psicoprofilaxis obstétrica, monitoreo fetal básico, planificación familiar y tamizaje de cáncer ginecológico. & Mañana / Tarde & 04 Lic. en Obstetricia \\
\addlinespace
Odontología & Odontología preventiva y restauradora, exodoncias simples, profilaxis dental y fluorización tópica. & Mañana / Tarde & 02 Cirujanos Dentistas \\
\addlinespace
Psicología & Atención de salud mental, tamizaje de violencia familiar, intervención en crisis y consejería psicológica. & Mañana & 02 Lic. en Psicología \\
\addlinespace
Tópico de Urgencias y Triaje & Evaluación de signos vitales, curaciones menores, suturas simples, administración de inyectables y nebulizaciones. & 12 Horas continuas & 04 Técnicos de Enfermería \\
\addlinespace
Farmacia SIS & Almacenamiento, custodia y dispensación ambulatoria de medicamentos del petitorio oficial para afiliados SIS. & 12 Horas continuas & 02 Técnicos de Farmacia \\
\addlinespace
Laboratorio Básico & Toma de muestras biológicas, hematocrito/hemoglobina, frotis sanguíneo, examen de orina y baciloscopías. & Mañana & 01 Tecnólogo / Técnico \\
\bottomrule
\end{tabularx}
\end{table}

El perfil epidemiológico de la demanda ambulatoria está hegemonizado por las Infecciones Respiratorias Agudas (IRA), Enfermedades Diarreicas Agudas (EDA), parasitosis intestinal, anemia por deficiencia de hierro en niños menores de 36 meses, infecciones del tracto urinario (ITU) en mujeres en edad fértil y gestantes, lumbalgias mecánicas derivadas del trabajo físico, y patologías crónicas no transmisibles (hipertensión arterial y diabetes mellitus tipo 2) en adultos mayores.

No obstante su relevancia estratégica, el establecimiento enfrenta una serie de \textbf{nudos críticos institucionales y barreras operativas} que inciden directamente en la subjetividad del usuario:
\begin{enumerate}
    \item \textbf{Sobrecarga asistencial y presión por rendimiento horario (tiempo de consulta abreviado):} El elevado número de usuarios asignados por turno y la exigencia de metas de producción asistencial reducen drásticamente los minutos dedicados a cada consulta clínica ambulatoria (frecuentemente de 10 a 15 minutos). Esta premura temporal condiciona un ritmo de atención apresurado que limita la profundización en el diálogo clínico, restringe el espacio para la consejería de enfermería y dificulta la acogida empática de las dudas y temores del paciente.
    \item \textbf{Limitaciones arquitectónicas y hacinamiento espacial:} El centro de salud funciona en un inmueble cuya infraestructura original ha sido sucesivamente adaptada, presentando pasillos estrechos, salas de espera techadas pero con escaso aislamiento térmico y bancas metálicas insuficientes. La tabiquería provisional entre consultorios adolece de aislamiento acústico, lo que provoca la filtración de conversaciones clínicas y quebranta el derecho a la privacidad e intimidad corporal y emocional de los pacientes.
    \item \textbf{Inestabilidad laboral y sobrecarga burocrática:} Una fracción considerable del personal profesional y técnico se encuentra bajo regímenes de contratación temporal transitoria (CAS), lo que origina una continua rotación de profesionales que interrumpe la continuidad del vínculo afectivo con los pacientes de la comunidad. A ello se suma la sobrecarga de digitación en los sistemas informáticos (HIS-MINSA y formato FUA-SIS), la cual absorbe gran parte del tiempo de consulta, obligando al profesional a interactuar con la pantalla del computador en detrimento del contacto visual con el usuario.
\end{enumerate}
